\documentclass{cvjc}
\usepackage[super]{nth}

\publicationlist{cvpub.bib}
\publicationlist{cvtalk.bib}

\begin{document}

\header{Joseph}{Clarke}

I am a Senior Scientist at the Met Office Hadley Centre. I am Earth System Scientist with particular interest in
nonlinear behaviour in the climate-carbon system. For example, I have
analysed Amazon dieback in \textsc{cmip6} data, used the complex
land-surface model \textsc{jules} to answer questions about
climate-carbon cycle feedbacks, and have designed conceptual models to
analyse rate induced tipping. I am interested in topics broadly
related to tipping points, such as overshoot, irreversibility,
feedbacks and sensitivity. As well as using models, I am interested in
empirical approaches to measuring changes in the stability of
potential tipping elements using early warning signals.  In addition
to my background in climate science, I use techniques from dynamical
systems, statistical and stochastic dynamics to understand these
nonlinear phenomena in the Earth system.

\section{Personal Information}
\begin{contact}
  \info{email}{\link{mailto:joseph.clarke@metoffice.gov.uk}{joseph.clarke@metoffice.gov.uk}}
  \info{website}{\link{https://josephjclarke.github.io/}{josephjclarke.github.io}}
  \info{orcid}{\link{https://orcid.org/0000-0002-6250-1543}{0000-0002-6250-1543}}
\end{contact}

\section{Education}
\begin{events}
  \event{2019--2023}{PhD in Mathematics, \textbf{University of Exeter},}{\textsc{PhD Thesis}: `Tipping Points and Early Warning Signals in the Climate-Carbon System', supervised by Professor Peter Cox and Professor Peter Ashwin.}
  \event{2015--2019}{MPhys Physics (\nth{1} Class), \textbf{Brasenose College, Oxford University},}{\textsc{Specialising in}: The Physics of Atmospheres and Oceans and Condensed Matter Physics \newline \textsc{Master's Thesis}: `Heat-engine based representation of convection with phase change in exoplanet atmospheres' with Professor Raymond Pierrehumbert.}
\end{events}

\section{Employment and Work Experience}
\begin{events}
  \event{2026--}{Senior Scientist, Met Office Hadley Centre}{\textsc{FUTURA} project. Researching abrupt shifts in the biosphere}
  \event{2023--2026}{Postdoctoral Research Associate, University of Exeter}{\textsc{predict} (Work Package lead) and \textsc{wiemip} projects. With \textsc{predict} I develop new early warning signals for tipping points. With \textsc{wiemip} I study warming induced emissions with a focus on overshoot and \textsc{tcre}.}
  \event{2019--2023}{Postgraduate Teaching Assistant, University of Exeter}{Assisted teaching in the following modules: `Probability, Statistics and Data', `Mathematics of Climate Change',
    `Ordinary Differential Equations', `Partial Differential Equations', `Fluid Dynamics'}
  \event{2018}{Summer Project with Dr Hannah Christensen, 
    University of Oxford.}{Took a dynamical systems approach to constraining climate
    sensitivity. I worked with Linear Inverse Models, Empirical Orthogonal Functions, 
    autoregressive models, Python,the Fluctuation Dissipation Theorem and undertook analysis of \textsc{cmip} data}
  \event{2017}{Summer Intern at the Earth Observation Data Group, University of Oxford.}{Developed a website to display observations of
    atmospheric anomalies in near real time. This involved processing the large amount of data generated
    and converting it to an understandable form.}
\end{events}

\section{Positions of Responsibility}
\begin{events}
  \event{2025--2026}{XCS Seminar Organiser}{I run the XCS (Exeter Climate Systems) seminars.}
  \event{2025}{Global Tipping Points Conference Committee}{Organising the Global Tipping Points 2025 conference, an interdisciplinary conference with over
  500 attendees.}
  \event{2024--2025}{GSI Seminar Organisation Committee}{We ran the Global Systems Institute seminar series. We invited speakers, organised rooms as well as
    provided AV support. Events included traditional seminars, assemblies and film screenings.}
  \event{2024}{Pint of Science Organiser}{Outreach activity arranging science talks for the public. Involved working with speakers, operating
    AV equipment and running the events.}
  % \event{2017--2019}{Web Officer, Physics Joint Consultative Committee (\textsc{pjcc})}{Responsible for maintaining the \textsc{pjcc} website using the Django framework. 
  %   Liaising with the physics department to represent student interests}
  % \event{2016}{Producer and Sound Technician, The Importance of Being Ernest,
  %   Brasenose Arts Week}{Responsible for operating microphones and playing music during performance. Designed and
  %   produced posters as well as acquired props. Assisted with rehearsals}
\end{events}


\section{Research Expertise}
\begin{itemize}
\item Scientific programming using Fortran and Python, including \textsc{gpu} programming using \textsc{jax}.
\item Running and developing land surface models
\item Climate Emulators
\item Emergent Constraints
\item Climate Sensitivity
\item Carbon Cycle biogeochemistry
\item Tipping Point Early Warning Signals
\item Autonomous and nonautonomous dynamical systems techniques (e.g. bifurcation analysis, overshoot, rate induced tipping, linear response theory)
\end{itemize}

\section{Publications}
\printpublications

\section{Selected Talks and Research Sessions}
\printtalks

\section{Successful Research Grant Applications}
\begin{events}
  \event{2025}{\textsc{predict}}{Funded by \textsc{esa}. Work Package lead}
  \event{2026}{\textsc{eo4tipp}}{Funded by \textsc{esa}}
\end{events}

\section{Research Supervision}
\begin{events}
  \event{2024-- }{Co-supervisor of PhD student Patric Boardman}{Thesis title: Emergent Constraints on Future Climate Extremes}
\end{events}
\section{Referees}
\begin{contact}
  \info{Professor Peter Cox}{\link{mailto:p.m.cox@exeter.ac.uk}{p.m.cox@exeter.ac.uk}}
  \info{Professor Peter Ashwin}{\link{mailto:p.ashwin@exeter.ac.uk}{p.ashwin.exeter.ac.uk}}
  % \info{Dr Paul Ritchie}{\link{mailto:paul.ritchie@exeter.ac.uk}{paul.ritchie.exeter.ac.uk}}
\end{contact}

\end{document}
